% Formulação atual do modelo MPPRP.
% Correspondência principal: src/solvers/MultProductProdctionRoutingProblem.py
% Separação dos cortes: src/solvers/RoundedCapacitySeparation.py

\documentclass[11pt,a4paper]{article}

\usepackage[T1]{fontenc}
\usepackage[utf8]{inputenc}
\usepackage[brazil]{babel}
\usepackage{amsmath,amssymb,mathtools}
\usepackage{booktabs}
\usepackage{geometry}
\usepackage{hyperref}
\usepackage{enumitem}

\geometry{margin=2.5cm}
\hypersetup{colorlinks=true,linkcolor=blue,urlcolor=blue}

\newcommand{\set}[1]{\mathcal{#1}}
\newcommand{\R}{\mathbb{R}}
\newcommand{\Zset}{\mathbb{Z}}
\newcommand{\ceil}[1]{\left\lceil #1 \right\rceil}
\newcommand{\pos}[1]{\left[#1\right]_+}

\title{Formulação Matemática Atual\\
       Problema Integrado de Produção--Roteamento Multiproduto}
\author{Implementação do repositório LSVRP}
\date{\today}

\begin{document}

\maketitle

\begin{abstract}
\sloppy
Este documento registra a formulação efetivamente implementada no modelo
\texttt{MultProduct\allowbreak{}Prodction\allowbreak{}Routing\allowbreak{}Problem}. O modelo combina dimensionamento
de lotes, estoques, entregas e roteamento de veículos em múltiplos períodos e
para múltiplos produtos. São incluídas as extensões de quebra de simetria HC1,
as desigualdades lógicas de Coelho e as desigualdades cumulativas de capacidade
arredondada, separadas heuristicamente em um \texttt{UserCutCallback} do CPLEX.
As duas primeiras extensões são restrições estáticas opcionais; a última é uma
família válida de cortes adicionada dinamicamente à relaxação corrente.
\end{abstract}

\tableofcontents

\section{Conjuntos e índices}

Sejam:
\begin{align*}
  \set{P} &= \{1,\ldots,P\} &&\text{produtos},\\
  \set{T} &= \{0,\ldots,T-1\} &&\text{períodos},\\
  \set{V} &= \{0,\ldots,V-1\} &&\text{veículos},\\
  \set{C} &= \{1,\ldots,N\} &&\text{clientes},\\
  \set{K} &= \{0\}\cup\set{C} &&\text{nós, com o depósito/planta }0,\\
  \set{J} &= \{1,2,3,4,5\} &&\text{componentes de custo}.
\end{align*}

No código, os índices de produtos, períodos, veículos e nós começam em zero.
Nesta formulação, a convenção matemática é usada para os clientes, enquanto o
nó $0$ continua representando a planta. Assim, um cliente $i\in\set{C}$ é o
nó $i$ nas variáveis de rota e corresponde à posição $i-1$ no vetor de demanda
da instância.

\section{Parâmetros}

\begin{description}[leftmargin=3.2cm,style=nextline]
  \item[$d_{pit}$] demanda do produto $p$ no cliente $i$ no período $t$;
  \item[$I^0_{pi}$] estoque inicial do produto $p$ no local $i$;
  \item[$b_p$] capacidade consumida para produzir uma unidade do produto $p$;
  \item[$B$] capacidade total de produção disponível em cada período;
  \item[$c_p$] custo unitário de produção do produto $p$;
  \item[$s_p$] custo de setup do produto $p$;
  \item[$M$] constante grande para a ativação da produção;
  \item[$U_{pi}$] capacidade máxima de estoque do produto $p$ no local $i$;
  \item[$C$] capacidade de carga de cada veículo;
  \item[$f$] custo fixo de transporte por saída da planta;
  \item[$a_{ik}$] custo variável do arco orientado $(i,k)$;
  \item[$M_p$] limite específico por produto, definido por
  $M_p=\min\{M,\sum_{\tau\in\set{T}}\sum_{i\in\set{C}}d_{pi\tau}\}$;
  \item[$D^{\mathrm{restante}}_{p,t}$] demanda do produto $p$ nos períodos
  $t,\ldots,T-1$;
  \item[$\overline Q_{pit}$] limite específico de entrega, definido por
  $\overline Q_{pit}=\min\{C,U_{pi}+d_{pit}\}$;
  \item[$w_j$] peso do componente de custo $j$;
  \item[$\alpha$] peso do termo minimax na função objetivo, com
  $0\le\alpha\le1$;
  \item[$b_t^j$] target original do componente $j$ no período $t$;
  \item[$\widetilde b_t^j$] target utilizado pelo modelo. Se o target original
  for zero, a implementação o substitui pela média temporal daquele componente.
\end{description}

\section{Configuração da formulação completa}

Para a variante fortalecida considerada neste documento, as extensões são
ativadas conjuntamente: modo multiobjetivo, HC1, desigualdades de Coelho e
Rounded Capacity Inequalities. Não há variáveis explícitas de visita. A opção
\texttt{solver.goal\_programming.positive\_only\_deviations} é verdadeira por
padrão, de modo que os desvios negativos $n_t^j$ não são criados.

O PSO não altera a formulação do MIP; quando ativado, ele fornece uma solução
inicial para o mesmo modelo exato.

\section{Variáveis de decisão}

\begin{align*}
  x_{pt} &\ge 0 &&\text{quantidade produzida de $p$ em $t$},\\
  y_{pt} &\in\{0,1\} &&\text{setup/ativação de $p$ em $t$},\\
  I_{pit} &\ge 0 &&\text{estoque de $p$ no local $i$ ao final de $t$},\\
  z_{vikt} &\in\{0,1\} &&\text{uso do arco orientado $(i,k)$ por $v$ em $t$},\\
  r_{pvikt} &\ge 0 &&\text{fluxo de $p$ no arco $(i,k)$ por $v$ em $t$},\\
  q_{pvit} &\ge 0 &&\text{entrega de $p$ pelo veículo $v$ ao cliente $i$ em $t$},\\
  p_t^j,n_t^j &\ge 0 &&\text{desvios positivo e negativo do componente $j$},\\
  \lambda &\ge 0 &&\text{maior desvio normalizado}.
\end{align*}

As variáveis de arco são criadas somente para $i\ne k$. As posições diagonais
não representam deslocamentos e não pertencem ao MIP.

\section{Componentes da função objetivo}

Os cinco componentes construídos no código são:
\begin{align}
  f_t^1 &= \sum_{p\in\set{P}} c_p x_{pt},
  &&\text{custo de produção}, \label{eq:f1}\\
  f_t^2 &= \sum_{p\in\set{P}} s_p y_{pt},
  &&\text{custo de setup}, \label{eq:f2}\\
  f_t^3 &= \sum_{p\in\set{P}}\sum_{i\in\set{K}} h_{pi}I_{pit},
  &&\text{custo de estoque}, \label{eq:f3}\\
  f_t^4 &= \sum_{v\in\set{V}}\sum_{k\in\set{C}} f z_{v0kt},
  &&\text{custo fixo de transporte}, \label{eq:f4}\\
  f_t^5 &= \sum_{v\in\set{V}}\sum_{i\in\set{K}}\sum_{k\in\set{K}:k\ne i}
           a_{ik}z_{vikt},
  &&\text{custo variável de transporte}. \label{eq:f5}
\end{align}

\section{Função objetivo multiobjetivo}

Quando \texttt{solver.multiobjective = true} e não está sendo construído o
target, o modelo resolve:
\begin{equation}
  \min\quad
  \alpha\lambda
  +(1-\alpha)\sum_{t\in\set{T}}\sum_{j\in\set{J}}
  \frac{w_j p_t^j}{\widetilde b_t^j}.
  \label{eq:objective}
\end{equation}

O vínculo do termo minimax com cada desvio é:
\begin{equation}
  w_jp_t^j \le \lambda\widetilde b_t^j
  \qquad \forall j\in\set{J},\ t\in\set{T}.
  \label{eq:lambda}
\end{equation}

Para a formulação padrão, os desvios são relacionados aos targets por:
\begin{equation}
  f_t^j+n_t^j-p_t^j=b_t^j
  \qquad \forall j\in\set{J},\ t\in\set{T}.
  \label{eq:deviation-balance}
\end{equation}

Opcionalmente, quando
\texttt{solver.goal\_programming.positive\_only\_deviations = true},
o código elimina $n_t^j$ e usa a projeção equivalente:
\begin{equation}
  p_t^j\ge f_t^j-b_t^j,
  \qquad p_t^j\ge0.
  \label{eq:positive-only}
\end{equation}

Essa é a configuração padrão atual. A variante com $n_t^j$ continua disponível
somente quando a opção é explicitamente definida como \texttt{false}.

\section{Restrições de produção e estoque}

\subsection{Balanço de estoque na planta}

Para cada produto e período, a produção mais o estoque anterior, menos as
entregas feitas a clientes, determina o estoque final na planta:
\begin{align}
  x_{p0}+I^0_{p0}-\sum_{v\in\set{V}}\sum_{i\in\set{C}}q_{pvi0}
    &= I_{p0,0}, && p\in\set{P}, \label{eq:plant-balance-initial}\\
  x_{pt}+I_{p0,t-1}-\sum_{v\in\set{V}}\sum_{i\in\set{C}}q_{pvit}
    &= I_{p0t}, && p\in\set{P},\ t\in\set{T}\setminus\{0\}.
  \label{eq:plant-balance}
\end{align}

Na primeira equação, o índice do estoque final é $I_{p0,0}$; na segunda,
o índice correto é $I_{p0t}$. A notação $I_{pit}$ usa $i=0$ para a planta.

\subsection{Balanço de estoque nos clientes}

Para cada cliente:
\begin{align}
  \sum_{v\in\set{V}}q_{pvi0}+I^0_{pi}-d_{pi0}
    &= I_{pi0}, && p\in\set{P},\ i\in\set{C}, \label{eq:customer-balance-initial}\\
  \sum_{v\in\set{V}}q_{pvit}+I_{pi,t-1}-d_{pit}
    &= I_{pit}, && p\in\set{P},\ i\in\set{C},\ t\in\set{T}\setminus\{0\}.
  \label{eq:customer-balance}
\end{align}

\subsection{Capacidade de produção, setup e estoque}

\begin{align}
  \sum_{p\in\set{P}} b_p x_{pt} &\le B,
  && t\in\set{T}, \label{eq:production-capacity}\\
  x_{pt} &\le U^X_{pt} y_{pt},
  && p\in\set{P},\ t\in\set{T}, \label{eq:setup-link}\\
  I_{pit} &\le U^I_{pit},
  && p\in\set{P},\ i\in\set{K},\ t\in\set{T}.
  \label{eq:inventory-capacity}
\end{align}

Os limites fortalecidos são:
\begin{align}
  U^X_{pt} &= \min\left\{M_p,\frac{B}{b_p},D^{\mathrm{restante}}_{p,t}\right\},
  \label{eq:production-upper-bound}\\
  M_p &= \min\left\{M,\sum_{\tau\in\set{T}}\sum_{i\in\set{C}}d_{pi\tau}\right\},
  \label{eq:product-big-m}\\
  D^{\mathrm{restante}}_{p,t} &= \sum_{\tau=t}^{T-1}\sum_{i\in\set{C}}d_{pi\tau}.
  \label{eq:remaining-demand}
\end{align}

Para a planta, $U^I_{p0t}$ é limitado pelo estoque inicial mais a produção
acumulada. Para um cliente $i\in\set{C}$:
\begin{equation}
  U^I_{pit}=\max\left\{0,\min\left\{U_{pi},I^0_{pi}+
  \sum_{s=0}^{t}\left(\overline Q_{pis}-d_{pis}\right)\right\}\right\},
  \qquad
  \overline Q_{pis}=\min\{C,U_{pi}+d_{pis}\}.
  \label{eq:inventory-upper-bound}
\end{equation}

Não há fator $V$ no limite dos clientes. A restrição de visita única limita a
soma das visitas entre veículos a um, e a conservação dos arcos iguala as
visitas de entrada e saída para cada veículo. Como $q_{pvit}$ é limitado por
$\overline Q_{pit}$ vezes a expressão de visita, a entrega total de cada
produto ao cliente em um período não excede $\overline Q_{pit}$. Somando o
balanço de estoque do cliente do período $0$ até $t$, subtrai-se a demanda
acumulada. O máximo com zero reflete a não negatividade de $I_{pit}$; o limite
da planta permanece inalterado.

\section{Fluxo de produtos e entregas}

Para cada cliente $k$, o fluxo de produto que entra menos o fluxo que sai é a
quantidade entregue naquele nó:
\begin{equation}
  \sum_{i\in\set{K}:i\ne k}r_{pvikt}
  -\sum_{l\in\set{K}:l\ne k}r_{pvklt}
  =q_{pvkt}
  \quad \forall p\in\set{P},\ v\in\set{V},\ k\in\set{C},\ t\in\set{T}.
  \label{eq:product-flow-customer}
\end{equation}

O balanço agregado na planta elimina fluxos de produto em subrotas
desconectadas:
\begin{equation}
  \sum_{v\in\set{V}}\sum_{k\in\set{C}}r_{pv0kt}
  -\sum_{v\in\set{V}}\sum_{i\in\set{C}}r_{pvi0t}
  =\sum_{v\in\set{V}}\sum_{l\in\set{C}}q_{pvlt}
  \quad \forall p\in\set{P},\ t\in\set{T}.
  \label{eq:product-flow-plant}
\end{equation}

\section{Roteamento de veículos}

\subsection{Capacidade nos arcos}

\begin{equation}
  \sum_{p\in\set{P}}r_{pvikt}\le C z_{vikt}
  \quad \forall v\in\set{V},\ i,k\in\set{K}:i\ne k,\ t\in\set{T}.
  \label{eq:arc-capacity}
\end{equation}

\subsection{Uso, conservação e visitação}

Cada veículo pode sair da planta no máximo uma vez em cada período:
\begin{equation}
  \sum_{k\in\set{C}}z_{v0kt}\le1
  \quad \forall v\in\set{V},\ t\in\set{T}.
  \label{eq:one-departure}
\end{equation}

A conservação de fluxo de veículos é imposta em todos os nós:
\begin{equation}
  \sum_{i\in\set{K}:i\ne k}z_{vikt}
  -\sum_{l\in\set{K}:l\ne k}z_{vklt}=0
  \quad \forall v\in\set{V},\ k\in\set{K},\ t\in\set{T}.
  \label{eq:vehicle-flow}
\end{equation}

Cada cliente pode ser visitado no máximo uma vez por período:
\begin{equation}
  \sum_{v\in\set{V}}\sum_{i\in\set{K}:i\ne k}z_{vikt}\le1
  \quad \forall k\in\set{C},\ t\in\set{T}.
  \label{eq:single-visit}
\end{equation}

\subsection{Bounds de entrega em função de $Z$}

Não é criada uma variável explícita de visita. Para cada veículo, cliente e
período, a expressão de visita é a soma dos arcos que saem do cliente:
\begin{equation}
  \operatorname{visit}_{vit}=\sum_{k\in\set{K}:k\ne i}z_{vikt}.
  \label{eq:z-visit-indicator}
\end{equation}

As entregas são limitadas diretamente por essa expressão:
\begin{align}
  \sum_{p\in\set{P}}q_{pvit}
  &\le C\sum_{k\in\set{K}:k\ne i}z_{vikt},
  &&v\in\set{V},\ i\in\set{C},\ t\in\set{T},
  \label{eq:z-visit-capacity}\\
  q_{pvit}
  &\le \overline Q_{pit}\sum_{k\in\set{K}:k\ne i}z_{vikt},
  &&p\in\set{P},\ v\in\set{V},\ i\in\set{C},\ t\in\set{T}.
  \label{eq:z-visit-product-capacity}
\end{align}

Essas restrições são padrão, e a conservação de fluxo mantém entrada e saída
iguais no cliente.

\section{Quebra de simetria HC1}

As restrições desta seção são adicionadas somente quando
\texttt{solver.symmetry\_breaking.hc1 = true}. Elas são válidas quando os
veículos são intercambiáveis, isto é, quando a troca de seus rótulos não muda
capacidade, custos, disponibilidade ou qualquer outra característica física.

Defina o indicador agregado de saída e o indicador de visita:
\begin{align}
  s_{vt} &= \sum_{k\in\set{C}}z_{v0kt},
  && v\in\set{V},\ t\in\set{T}, \label{eq:vehicle-active}\\
  u_{vit} &= \sum_{k\in\set{K}:k\ne i}z_{vikt},
  && v\in\set{V},\ i\in\set{C},\ t\in\set{T}.
  \label{eq:customer-visit}
\end{align}

Embora $s$ e $u$ sejam apenas expressões auxiliares e não variáveis criadas
no MIP, eles tornam as restrições implementadas explícitas.

\subsection{VC: ordenação da ativação dos veículos}

\begin{equation}
  s_{vt}\le s_{v-1,t}
  \quad \forall v\in\set{V}\setminus\{0\},\ t\in\set{T}.
  \label{eq:VC}
\end{equation}

\subsection{HC1: ordenação das visitas}

\begin{equation}
  u_{vit}\le \sum_{j\in\set{C}:j<i}u_{v-1,j,t}
  \quad \forall v\in\set{V}\setminus\{0\},\ i\in\set{C},\ t\in\set{T}.
  \label{eq:HC1}
\end{equation}

Para $i=1$, o lado direito é uma soma vazia e vale zero. Portanto, um veículo
de índice maior não pode ser o primeiro veículo a atender o primeiro cliente.

\section{Desigualdades lógicas de Coelho}

As desigualdades de Coelho são adicionadas quando
\texttt{solver.coelho\_inequalities = true}. O artigo original usa uma aresta
não orientada $x$ e uma variável de visita $y$; a implementação atual não cria
$y$ e projeta essas relações sobre os arcos orientados $z$ por meio de $s$ e
$u$ definidos em \eqref{eq:vehicle-active}--\eqref{eq:customer-visit}.

\subsection{Coelho (15): arco planta--cliente}

\begin{equation}
  z_{v0it}\le 2u_{vit}
  \quad \forall v\in\set{V},\ i\in\set{C},\ t\in\set{T}.
  \label{eq:coelho-15}
\end{equation}

\subsection{Coelho (16): arco implica visita da origem}

Para arcos que saem da planta:
\begin{equation}
  z_{v0kt}\le s_{vt}
  \quad \forall v\in\set{V},\ k\in\set{C},\ t\in\set{T}.
  \label{eq:coelho-16-plant}
\end{equation}

Para arcos que saem de um cliente:
\begin{equation}
  z_{vikt}\le u_{vit}
  \quad \forall v\in\set{V},\ i\in\set{C},\ k\in\set{K}:k\ne i,\ t\in\set{T}.
  \label{eq:coelho-16-customer}
\end{equation}

\subsection{Coelho (17): visita implica uso do veículo}

\begin{equation}
  u_{vit}\le s_{vt}
  \quad \forall v\in\set{V},\ i\in\set{C},\ t\in\set{T}.
  \label{eq:coelho-17}
\end{equation}

\section{Desigualdades cumulativas de capacidade arredondada}
\label{sec:rounded-capacity}

\subsection{Demanda obrigatória}

Para um prefixo temporal $0,\ldots,\tau$, a demanda obrigatória do produto
$p$ no cliente $i$ é definida por:
\begin{equation}
  d^{\mathrm{obrig}}_{pi}(\tau)
  =\pos{\sum_{t=0}^{\tau}d_{pit}-I^0_{pi}}.
  \label{eq:obligatory-demand-product}
\end{equation}

O volume agregado do subconjunto $S\subseteq\set{C}$ é:
\begin{equation}
  D(S,\tau)=\sum_{i\in S}\sum_{p\in\set{P}}
  d^{\mathrm{obrig}}_{pi}(\tau).
  \label{eq:obligatory-demand-subset}
\end{equation}

O estoque inicial é descontado separadamente por produto antes da agregação.
Assim, o estoque de um produto não pode compensar a demanda de outro.

\subsection{Arcos que cruzam um subconjunto}

Para $S\subseteq\set{C}$, defina:
\begin{equation}
  \delta(S,t)=
  \sum_{v\in\set{V}}
  \sum_{i\in\set{K}}\sum_{k\in\set{K}:k\ne i}
  \mathbf{1}\!\left[
    \mathbf{1}[i\in S]\ne\mathbf{1}[k\in S]
  \right]z_{vikt}.
  \label{eq:cut-crossing}
\end{equation}

Como os arcos são orientados, as duas orientações que cruzam o corte são
somadas. Para uma rota fechada, uma travessia de entrada e uma de saída
produzem o fator $2$ da desigualdade seguinte.

\subsection{Família válida de cortes}

Para todo subconjunto de clientes $S$ e todo prefixo $\tau$:
\begin{equation}
  \sum_{t=0}^{\tau}\delta(S,t)
  \ge
  2\ceil{\frac{D(S,\tau)}{C}}.
  \label{eq:rounded-capacity}
\end{equation}

Essa desigualdade expressa que, para atender o volume obrigatório de $S$, são
necessários pelo menos $\ceil{D(S,\tau)/C}$ carregamentos de capacidade $C$;
cada carregamento precisa cruzar o corte para entrar e sair de $S$.

\subsection{Separação implementada}

As desigualdades em \eqref{eq:rounded-capacity} não são todas inseridas na
construção inicial do modelo. Quando o recurso está habilitado, o CPLEX chama
um \texttt{UserCutCallback} em pontos fracionários da relaxação. A rotina:

\begin{enumerate}[leftmargin=1.8em]
  \item lê os valores fracionários das variáveis $z$;
  \item agrega os arcos por prefixo temporal;
  \item cria sementes unitárias e pares de clientes com maior conectividade
        interna;
  \item aplica crescimento e uma busca local vetorizada por remoção, adição e
        troca de clientes;
  \item calcula a violação de cada candidato e retém os mais violados;
  \item elimina cortes duplicados por meio da chave $(\tau,S)$;
  \item adiciona as linhas globais com sentido $\ge$ quando a violação supera
        \texttt{min\_violation}.
\end{enumerate}

O separador é heurístico: ele não enumera todos os subconjuntos $S$. A família
em \eqref{eq:rounded-capacity} é válida; a heurística somente determina quais
de suas desigualdades serão encontradas em cada chamada.

\subsection{Agendamento e limites}

Os parâmetros atuais são:
\begin{center}
\begin{tabular}{ll}
\toprule
Parâmetro & Valor padrão atual \\
\midrule
\texttt{enabled} & \texttt{false} \\
\texttt{node\_frequency} & $50$ \\
\texttt{max\_non\_root\_node} & $200$ \\
\texttt{max\_cuts\_per\_callback} & $20$ \\
\texttt{max\_cuts\_per\_non\_root\_callback} & $3$ \\
\texttt{min\_violation} & $10^{-6}$ \\
\bottomrule
\end{tabular}
\end{center}

O callback é sempre elegível no nó raiz. Nos nós não-raiz, ele só é elegível
até o nó configurado e em múltiplos de \texttt{node\_frequency}. O limite de
20 é por chamada do callback; se o CPLEX chamar a separação novamente, novas
desigualdades ainda podem ser adicionadas, sujeitas à deduplicação.

Os cortes são instalados somente no MIP exato. Eles não são adicionados ao
PSO, ao modelo auxiliar de dimensionamento de lotes
(\texttt{LotSizingRelaxation}) nem ao modelo usado para construir targets.
Também não é criado um segundo CPLEX dentro do callback.

\section{Domínios e formulação completa}

Além das restrições anteriores, valem:
\begin{align}
  x_{pt},I_{pit},r_{pvikt},q_{pvit},p_t^j,\lambda &\ge0,
  \label{eq:nonnegative}\\
  y_{pt},z_{vikt} &\in\{0,1\}.
  \label{eq:binary}
\end{align}

O modelo atual é, portanto, um MIP formado por
\eqref{eq:objective}--\eqref{eq:lambda}, pelos balanços de estoque
\eqref{eq:plant-balance-initial}--\eqref{eq:customer-balance}, pelas
capacidades \eqref{eq:production-capacity}--\eqref{eq:inventory-capacity},
pelos fluxos \eqref{eq:product-flow-customer}--\eqref{eq:product-flow-plant},
pelas restrições de rota \eqref{eq:arc-capacity}--\eqref{eq:single-visit},
pelos bounds de entrega \eqref{eq:z-visit-capacity}--
\eqref{eq:z-visit-product-capacity}, por HC1 em \eqref{eq:VC}--\eqref{eq:HC1} e Coelho em
\eqref{eq:coelho-15}--\eqref{eq:coelho-17}, além dos cortes dinâmicos da
família \eqref{eq:rounded-capacity}.

\section{Correspondência com a implementação}

\begin{center}
\begin{tabular}{ll}
\toprule
Componente & Implementação \\
\midrule
Variáveis e restrições base & \texttt{MultProductProdctionRoutingProblem.py} \\
Função objetivo & \texttt{crateObjectiveFunction} \\
Balanços de estoque & \texttt{createEstablishInvetoryBalanceAtPlant},\\
& \texttt{creteInventoryBalancingInventoryCustomers} \\
Fluxo de produtos & \texttt{createVehiclePreventTransshipmentIntermediateNodes},\\
& \texttt{createEliminationSubroutes} \\
Capacidade e roteamento & métodos \texttt{createVehicle...} \\
Bounds de entrega em função de $Z$ & \texttt{createVehicleVisitDeliveryBounds} \\
HC1/VC & \texttt{createVehicleSymmetryBreaking} \\
Coelho & \texttt{createCoelhoValidInequalities} \\
Separador & \texttt{RoundedCapacitySeparation.py} \\
Callback & \texttt{\_install\_rounded\_capacity\_callback} \\
\bottomrule
\end{tabular}
\end{center}

\end{document}
